\documentclass[letterpaper,11pt]{article}
\usepackage[T1]{fontenc}
\usepackage[utf8]{inputenc}
\usepackage[hidelinks]{hyperref}
\usepackage{xcolor}
\usepackage[margin=1in]{geometry}
\usepackage[default]{sourcesanspro}
\definecolor{ResumeNavy}{HTML}{18324A}
\definecolor{ResumeSlate}{HTML}{5B6573}
\color{black}
\setlength{\parskip}{8pt}
\setlength{\parindent}{0pt}
\hypersetup{colorlinks=true, urlcolor=ResumeNavy}
\begin{document}
\begin{center}
  {\LARGE \bfseries \textcolor{ResumeNavy}{Deva Anand}} \\[3pt]
  {\small
  \href{mailto:devaanand@umass.edu}{devaanand@umass.edu} \enspace\textbar\enspace
  (413) 315-1715 \enspace\textbar\enspace
  \href{https://github.com/Deva-1903}{github.com/Deva-1903} \enspace\textbar\enspace
  \href{https://linkedin.com/in/devaaa/}{linkedin.com/in/devaaa}}
\end{center}
\vspace{6pt}

June 10, 2026

\vspace{2pt}
DataVisor \\
AI/ML Engineering Internship Hiring Team

\vspace{6pt}
Dear Hiring Team,

I'm applying for the AI/ML Engineering Intern role. I'm a Computer Science master's student at UMass Amherst, and what stands out about this position is the blend it asks for, data pipelines, ML, agentic AI, and a real concern for security, which happens to be the exact mix of things I've spent the last couple of years working on.

The data-pipeline side is where I'm most at home. I built a Spark pipeline that processed a 1.4-billion-row dataset and cut its end-to-end runtime by about 25\% by reworking joins, enabling adaptive query execution, and tuning partitioning, then used stage-level profiling to bring the analytics tail-latency ratio down meaningfully. That experience of making a high-throughput pipeline both faster and more predictable is directly what your Intelligence Layer work calls for.

On the AI side, most of my recent projects are agentic. I built \textbf{AgenticSearch}, an LLM agent that plans retrieval, uses tools, and verifies its own outputs with cell-level provenance, hardened with 150+ tests; \textbf{AutoEval}, a self-improving evaluation agent; and a team \textbf{RAG} pipeline where my multi-view retrieval lifted supporting-fact recall by 2.68\%. I've also trained supervised models from scratch (CNN/ResNet classifiers in CS~689) and published a transfer-learning classifier, so the ``deploy and monitor ML pipelines'' side isn't unfamiliar either.

Finally, security isn't an afterthought for me. At Wysa I owned reliability and security for a multi-tenant healthcare platform, implementing AWS KMS envelope encryption for sensitive data and remediating 20+ penetration-testing findings, which is why a fraud-and-risk platform built around privacy-first design genuinely appeals to me.

I know I work with Spark rather than Kafka or Flink today, but those are the kinds of tools I pick up quickly, and I'd be excited to learn them on real production systems. I'd welcome the chance to contribute.

\vspace{4pt}
Thank you for your time and consideration.

\vspace{8pt}
Sincerely, \\
Deva Anand
\end{document}
